% Chapter 3.8: Synthesis Matrix and Real-World Deployment Roadblocks
% Part III of the book.

\chapter{Synthesis Matrix and Real-World Deployment Roadblocks}
\label{ch:part3-synthesis-matrix-and-real-world-deployment}
\label{sec:synthesis}
\index{synthesis matrix|textbf}\index{deployment roadblocks|textbf}
% -----------------------------------------------------------------------------

The five application chapters of this part each replaced one solver with one surrogate, and
the ledger format this Part uses was meant to report two numbers for each: a ratio of wall-clock times and an
error against a reference field. Three of the five get both; the cardiac and tomography rows
carry neither a time ratio nor a measured field error, and the matrix says so. Those
numbers are defensible one at a time; taken together --- the way a design office or a clinic must
take them --- they are not comparable. The aerodynamic multiplier was formed
against a volume discretization of computational fluid dynamics (CFD) on a graphics processing
unit (GPU); the reservoir multiplier against an implicit finite difference simulator on a
central processing unit (CPU); the atmospheric multiplier against the spectral-transform
dynamics of the European Centre for Medium-Range Weather Forecasts (ECMWF) Integrated
Forecasting System (IFS); and in the tomographic row the comparison is not a timing but
a reconstruction quality at fixed radiation dose.\index{speedup factor!provenance of} A
synthesis records provenance before it records figures.
Table~\ref{tab:synthesis_matrix} does that, and adds the column the inherited accounting
leaves out:
which rigorous statement, if any, Part~I supplies for the architecture in the row.

The rest of the chapter concerns what a multiplier does not price. Every application here
was chosen because the query count $Q$ of \eqref{eq:ch01-cost-bookkeeping} --- $M$ as Part~I writes it --- is large, so the
amortization argument of
\chapref{ch:part1-introduction-the-paradigm-shift-to-infinit} holds; deployment is governed
by obstacles that argument does not reach. A certifier cannot run the reference solver on
the field being decided upon, so a reported error is anchored to a machine no longer in the
loop. A conservation law that a finite volume method (FVM) satisfies cell by cell is
satisfied by a trained network only on average. The labels that made the surrogate possible
were themselves solver runs, and that cost recurs with every retraining. A guarantee stated
over a compact subset of the training measure --- the hypothesis
Assumption~\ref{ass:ch09-compact} makes explicit --- licenses nothing about the geometry,
aquifer, pacing protocol, or scanner that arrives next.

Prerequisites are the book's: the error split \eqref{eq:ch01-error-split} in ledger
form \eqref{eq:ch06-decomp}, the architecture comparison of
Tables~\ref{tab:comparison} and~\ref{tab:ch8-ledger}, the deployment budget
\eqref{eq:ch27-budget}, and the five application chapters. Section~\ref{sec:ch38-matrix}
builds the matrix; Section~\ref{sec:ch38-columns} reads its columns against the arithmetic
that produced them;
Sections~\ref{sec:ch38-roadblocks} and~\ref{sec:ch38-data-shift} state the obstacles, each
with the estimate that makes it precise rather than rhetorical; Section~\ref{sec:ch38-closing}
returns to the question the book opened with and hands off to the back matter. No new
measurement, architecture, or theorem appears here; a closing chapter's value is that its
claims are already in print and can be held against one another.

% -----------------------------------------------------------------------------
\section{The Synthesis Matrix and What Its Columns Mean}
\label{sec:ch38-matrix}
% -----------------------------------------------------------------------------

Table~\ref{tab:synthesis_matrix} is ordered by application rather than by architecture,
because the domain fixes the equation, the mesh, the metric, and the reviewer. Each row is
the version of the multi-query surrogate problem \eqref{eq:ch01-solution-operator} that its
chapter posed: the Reynolds-averaged Navier--Stokes (RANS) equations for a rigid body in a
free stream, the hydrostatic primitive equations on a rotating sphere, two-phase Darcy flow
in a heterogeneous aquifer, the monodomain reaction--diffusion system on a patient-specific
ventricle, and the inverse Radon transform under angular
undersampling.\index{Darcy flow!two-phase}\index{monodomain equation}\index{Radon transform!inverse problem}\index{Reynolds-averaged Navier--Stokes (RANS) equations}\index{hydrostatic primitive equations}
The five rows are served by five architectures: the geometry-informed neural operator (GINO)
\cite{li2023geometry}, FourCastNet \cite{pathak2022fourcastnet} built on the adaptive Fourier
neural operator (AFNO), the U-Net-enhanced Fourier neural operator (U-FNO)
\cite{wen2022ufno}, the physics-informed neural operator (PINO) \cite{li2021physics}
applied to the monodomain system with an Aliev--Panfilov ionic gate
\cite{aliev1996simple}, and, in the tomographic row, a reconstruction built on the
Fourier neural operator architecture \cite{li2021fourier} as its layer --- the cite
names the layer, no published reconstruction system --- arranged either as an unrolled data-consistency
cascade or as a learned filtering stage inside filtered backprojection (FBP). The first
three are named models with published evaluations; the last two are formulations this
book's own chapters build rather than released networks, which is why their rows are
scored on quantities rather than on a published model's figures. Those quantities are the
drag coefficient $C_d$, the anomaly
correlation coefficient (ACC), the plume-restricted saturation mean-absolute error and the
relative pressure-buildup error, and the
image-quality measures defined in \eqref{eq:ch36-metrics}: peak signal-to-noise ratio
(PSNR), structural similarity index measure (SSIM), and root-mean-square error (RMSE) in
Hounsfield units (HU).

\begin{table}[t]
\centering
\footnotesize
\caption{\textbf{Synthesis matrix: application, governing system, architecture, cost ratio as formed, fidelity as reported, and the rigorous statement Part~I supplies.} Each of the three printed ratios has a different provenance: the aerodynamic row prints the cost--accuracy ratio reported by \cite{li2023geometry}; the weather row divides the study's sub-two-second surrogate latency, carried here as $1.8$\,s, by a $120$-minute per-member IFS figure this Part attributes (the study times whole ensembles: $82$\,min of wall-clock for a $15$-day, $51$-member run, Table~\ref{tab:benchmarks_hardware}); the carbon-storage row divides the $10$\,min average solve reported by \cite{wen2022ufno} by a query cost of $0.02$\,s --- the published $0.01$\,s per-field time doubled, because saturation and pressure are returned by two networks, as
\chapref{ch:part3-subsurface-multiphase-flow-and-carbon-stor} records. Every ratio is in essence a hardware-bound system comparison rather than an algorithmic one. Where no published pair of timings exists the ratio cell says so instead of printing a multiplier; a fidelity cell carries a measured claim only where a published study states one, and otherwise records this book's own accounting in the quantities its chapter defines. The final column uses the convention of Table~\ref{tab:ch8-ledger}: ``none here'' records that this book develops or cites no theorem for that architecture, not that the literature has no result.}
\label{tab:synthesis_matrix}
\begin{tabularx}{\textwidth}{@{}l >{\raggedright\arraybackslash}X l l >{\raggedright\arraybackslash}X >{\raggedright\arraybackslash}X@{}}
\toprule
\textbf{Domain} & \textbf{Governing system} & \textbf{Operator} & \textbf{Cost ratio} & \textbf{Fidelity as reported} & \textbf{Guarantee in Part~I} \\
\midrule
Aerodynamics
& steady incompressible RANS
& GINO \cite{li2023geometry}
& $2.6\times10^{4}$
& $C_d$ within $3\%$ at the reported cost ratio
& none here: neither graph stage nor latent grid proved
\\
\addlinespace
Global weather
& hydrostatic primitive equations on $\mathbb{S}^{2}$
& AFNO \cite{pathak2022fourcastnet}
& $4\times10^{3}$
& parity with IFS to day~3; $\mathrm{ACC}>0.6$ for TCWV beyond day~8
& none here: the rate theorem holds on the torus
\\
\addlinespace
Carbon storage
& two-phase Darcy flow, $S_w+S_g=1$
& U-FNO \cite{wen2022ufno}
& $3\times10^{4}$
& plume saturation MAE $1.6\%$, pressure error $0.68\%$
& truncation tail at a front, \eqref{eq:ch09-jump-tail}
\\
\addlinespace
Cardiac
& mono\-domain reaction--diffusion, anisotropic $\mathbf{D}(x)$
& PINO \cite{li2021physics}
& not reported
& no measured field error here; residual term \eqref{eq:ch07-pde-loss} only
& no parabolic analogue of \eqref{eq:ch07-apriori}
\\
\addlinespace
Tomography
& inverse Radon transform, angularly undersampled
& FNO-class, \chapref{ch:part3-medical-imaging-and-tomographic-reconstruc}
& not established
& no operating point carried; metric conventions \eqref{eq:ch36-metrics}
& none here: no stability constant for the inverse
\\
\bottomrule
\end{tabularx}
\end{table}

Read the cost column first: it is where a synthesis is sloppiest. Each printed
entry is either the aerodynamic cost--accuracy ratio, a published latency divided by
a baseline this Part attributes, or a published solve time divided by a per-field query
cost doubled by this Part's own two-network accounting --- the
provenances this matrix's own caption lays out, of which Table~\ref{tab:benchmarks_hardware}'s caption shows two and adds a further kind, book-assembled estimates. The
baseline each entry is priced against is a choice rather than a
fact. The baselines are recoverable from this part's own chapters: an OpenFOAM-class finite
volume solve for aerodynamics, an Eclipse-class implicit finite difference simulator for the
reservoir, openCARP finite elements for the cardiac row, the IFS for the atmosphere, and FBP
for tomography. The first of those names is already doing work the aerodynamic row's multiplier cannot support. The
$2.6\times10^{4}$ of the aerodynamic row is the cost--accuracy ratio at $3\%$ drag error
reported by \cite{li2023geometry}, formed against the study's own GPU-accelerated
volume-CFD solve --- the multi-core, two-V100 label-generation configuration of
\chapref{ch:part3-external-aerodynamics-and-surface-flows}, whose wall-clock time the study
does not pair with the ratio --- rather than against the CPU alternative quoted in that
chapter: $15$--$45$ minutes on a many-core node, itself flagged there as an
order-of-magnitude pairing this Part inherits, not a time the study reports for that
configuration; price the same GINO surrogate against that alternative and the multiplier
changes again. Neither quotient describes an algorithm alone: the
printed one pits a multi-GPU solve against a single-GPU evaluation, the alternative an
accelerator against a CPU cluster.\index{wall-clock ratio!choice of baseline}
The two cells this part's inherited ledger filled with $1{,}500\times$ and $100\times$
are entered here as unreported or unestablished: neither is traceable to a pair of timings over the
baseline its row names. The cardiac row carries no multiplier because the comparison
there is against a fine-mesh finite element solve whose cost this book does not price;
the tomographic row carries none because what is bought there is image quality at a
dose rather than time, and the reconstruction's expense is reckoned in reprojection
passes \eqref{eq:ch36-consistency}, not in seconds saved. Claims of parity with FBP's
runtime, or of a wide margin over other learned reconstructions, would be claims about
particular reconstructors, and this book makes none. An unreconstructable number in a table is how an
unsupported figure acquires a citation, and a synthesis is where the circulation stops.

% -----------------------------------------------------------------------------
\section{The Cost Column's Second Reading, Fidelity, and the Guarantee Column}
\label{sec:ch38-columns}
% -----------------------------------------------------------------------------

The cost column has a second reading. Expand the training cost into what it is --- solver runs
plus an optimization --- and the crossover query count of
\eqref{eq:ch01-cost-bookkeeping} becomes
\begin{equation}
  C_{\mathrm{sur}}(Q)
  \;=\;
  \underbrace{n_{\mathrm{lab}}\,c_{\mathrm{sim}}(N_{\mathrm{lab}}) + c_{\mathrm{opt}}}_{c_{\mathrm{train}}}
  + Q\,c_{\mathrm{eval}}(N),
  \qquad
  Q^{\star}
  \;=\;
  \frac{n_{\mathrm{lab}}\,c_{\mathrm{sim}}(N_{\mathrm{lab}}) + c_{\mathrm{opt}}}
       {c_{\mathrm{sim}}(N) - c_{\mathrm{eval}}(N)}.
  \label{eq:ch38-crossover}
\end{equation}
\index{crossover query count}
Equation \eqref{eq:ch38-crossover} is arithmetic; its consequence is a design rule. Where
$c_{\mathrm{eval}}(N)<c_{\mathrm{sim}}(N)$ and the labels were generated on a grid at least
as fine as the one served, $c_{\mathrm{sim}}(N_{\mathrm{lab}})\ge c_{\mathrm{sim}}(N)$ pushes
$Q^{\star}$ above the number of labeled solves $n_{\mathrm{lab}}$. A surrogate must answer
more fields than it was trained on before it costs less than the solver it replaced, and
labeling more finely than one serves raises the threshold without any change in the network.
That is why the five rows were worth building. Sampling a permeability or stiffness field by
Markov chain Monte Carlo (MCMC) is stated in
\chapref{ch:part3-the-multi-query-bottleneck-in-computationa} to want $10^{5}$--$10^{7}$
forward evaluations; a fifty-member atmospheric ensemble run twice daily is $36{,}500$ solves
in a year; an ablation plan is queried per candidate lesion under a latency no tetrahedral
solve meets. A single verification run should by contrast stay with a classical
solver, the qualification attached to \eqref{eq:ch01-cost-bookkeeping} in
\chapref{ch:part1-introduction-the-paradigm-shift-to-infinit}. The tomographic row differs in
kind: the operator is bought for image quality at lower dose rather than for wall-clock, and
the two defects it is bought against are not interchangeable. Sparse-view acquisition leaves
the angular arc complete with too few samples inside it, so one trained model may be asked
to serve several projection sampling rates; the limited-angle case removes an arc of the
acquisition entirely, and no density of sampling inside what remains recovers it.
\chapref{ch:part3-medical-imaging-and-tomographic-reconstruc} separates the two
arrangements along exactly that line. The gain figures an earlier version of this table
carried --- a margin at $18$ views, and an average margin over sampling rates measured
against convolutional baselines --- are dropped rather than restated: they were measured on
reconstructors this book does not name, and a point on a curve and an average over rates
are in any case different objects, which is why neither was convertible into the other.

The fidelity column asks for the same skepticism. Every metric printed
there is computable, which is why it was reported; none is the quantity a decision reads. A
$3\%$ error in $C_d$ is an integral over the surface, constraining neither the separation
location nor the wall shear stress (WSS) a drag-reduction change actually moves. The ACC is a
spatial correlation of two fields at one lead time, and \eqref{eq:ch09-weak} is this book's
reason for refusing to extrapolate it to the long-horizon statistics a climate application
consumes. A $1.6\%$ mean absolute saturation error inside the plume says nothing about where the
$\mathrm{CO}_2$ front sits or how much flux crosses the caprock, and the Gibbs analysis
around \eqref{eq:ch09-gibbs} explains why the number stays quiet: under a volume
average, whether in $L^2$ or in $L^1$, an error of fixed height on
a set of vanishing measure costs little in either average and a great deal in a permit. PSNR, SSIM,
and RMSE in HU \eqref{eq:ch36-metrics} are the most decision-adjacent of the five, being defined against a reference
reconstruction a radiologist accepts --- a reference that exists in the evaluation and not in
the clinic, where only the undersampled sinogram is
available.\index{fidelity metrics!as proxies for decisions}\index{tomography!sparse-view reconstruction}\index{aerodynamics!surface quantities not priced by drag error}\index{weather forecasting!skill over long horizons}

The last column is the one this book can sign. None of the five architectures above
carries a convergence theorem proved for it as used: Table~\ref{tab:ch8-ledger} gives GINO a
row and no theorem; the spectral layer behind FourCastNet is not the
translation-invariant multiplier of Theorem~\ref{thm:fno_convergence}, whose hypotheses live
on the torus; U-FNO contributes exactly the multi-scale content that rate does not price, and
\eqref{eq:ch09-jump-tail} shows what a front costs it; a reconstruction whose kernels are
parameterized on the continuum decouples its weights from the detector sampling density but
inherits no modulus of continuity for an undersampled inverse problem; and the arrangement
that keeps the analytic ramp and backprojection \eqref{eq:ch36-fbp} intact leaves the learned
part to extrapolate the missing angular wedge, for which this book supplies no
analysis.\index{Geometry-Informed Neural Operator (GINO)!in deployment}\index{adaptive Fourier neural operator (AFNO)}\index{U-Net-enhanced Fourier neural operator (U-FNO)!multi-scale content}\index{tomography!operator-learning reconstruction}\index{Physics-Informed Neural Operator (PINO)!cardiac application}
PINO is the one row with a theorem adjacent to it, and the adjacency is conditional: the
residual-to-error implication \eqref{eq:ch07-apriori} is proved there for a coercive
divergence-form problem whose trial field already has the correct trace, and no parabolic
counterpart is derived in this book for the monodomain system with its nonlinear ionic
current. The matrix thus records an exchange rate of five to one: real-time and lower-dose
answers, purchased with guarantees never attached to the purchased object.

% -----------------------------------------------------------------------------
\section{Translational Roadblocks to Production Deployment}
\label{sec:ch38-roadblocks}
% -----------------------------------------------------------------------------

The first obstacle is the one the ratio conceals: error estimation without a reference
solver. In the classical workflow a reviewer asks for a residual, a refinement study, or
both, because the solver is available and accuracy is purchasable at query
time --- the point \chapref{ch:part2-production-deployment-and-high-throughput-} makes against
its own speedup table. A deployed surrogate has no such dial: mode counts, width, and depth
are frozen in the archive of Implementation Pattern~\ref{imp:ch27-artifact}, so more compute
per query buys no accuracy. The principled substitute is a stability estimate, bounding the
error in the functional actually decided upon by the residual the surrogate's own output
produces:
\begin{equation}
  \big|J(u)-J(v)\big|
  \;\le\;
  L_J\,\big\|v-u\big\|_{\mathcal{U}}
  \;\le\;
  L_J\,C_{\mathrm{st}}\,\big\|\mathcal{N}(v;a)\big\|_{X^{*}(\Omega)},
  \qquad v = \mathcal{G}_{\theta}(a),
  \label{eq:ch38-residual}
\end{equation}
\index{error estimation!without a reference solver}\index{certification of learned surrogates}
Here $J$ is a functional on the output space $\mathcal{U}$, Lipschitz with constant $L_J$;
$\mathcal{N}(\,\cdot\,;a)$ is the residual operator of \eqref{eq:ch07-pde-loss}; $X^{*}(\Omega)$
is the dual test space measuring the weak residual; $u$ is the exact state; and
$C_{\mathrm{st}}$ is the stability constant of the underlying a priori estimate --- for Darcy
flow under the uniform ellipticity of Assumption~\ref{ass:ellipticity},
$C_{\mathrm{st}}=1/a_{\min}$ in the energy norm. Equation \eqref{eq:ch38-residual} is the
deployment form of \eqref{eq:ch07-apriori} and inherits its conditions: the trial field must
carry the correct trace, coercivity or an equivalent inf-sup condition must hold, and the
bound degrades as the ellipticity constant does. Read against the matrix, that machinery
exists in model form for at most one row --- the elliptic Darcy parent of the subsurface
problem --- and for none of the five as deployed. It has one further merit and one further
limit: the residual in \eqref{eq:ch38-residual} is computable at inference from the output
and the equation alone, with no label required, but its cost and its fidelity depend on being
able to differentiate the surrogate's output field --- exactly on the uniform periodic grids
of \chapref{ch:part2-fourier-neural-operators-fno-from-scratch}, approximately and at added
expense on the surface stages of \chapref{ch:part3-external-aerodynamics-and-surface-flows}. Where the estimate is
unavailable, the substitute is a schedule rather than a theorem: run the reference
solver on a randomly selected fraction of queries. That returns the audit to
\eqref{eq:ch38-crossover} as a term $Q_{\mathrm{audit}}\,c_{\mathrm{sim}}(N)$ and preserves
the surrogate's advantage exactly while $Q_{\mathrm{audit}}\ll Q$. Certification by sampling
is slower than inference and quicker than hope.

The second obstacle is conservation. Certifiers ask for guaranteed balance of mass,
momentum, and energy; an $L^2$-trained network delivers a bound on a global quantity and
nothing more. Let $\mathcal{T}(w)=\int_{\Omega}w\,\mathrm{d}x$ be the conserved total of a balance law
with no flux through $\partial\Omega$. H{\"o}lder's inequality converts the training error
into a defect estimate,
\begin{equation}
  \Delta_{\mathcal{T}}(v,u)
  \;:=\;
  \frac{\big|\mathcal{T}(v)-\mathcal{T}(u)\big|}{\big|\mathcal{T}(u)\big|}
  \;\le\;
  \frac{1}{\big|\mathcal{T}(u)\big|}\int_{\Omega}\!\big|v-u\big|\,\mathrm{d}x
  \;\le\;
  \frac{|\Omega|^{1/2}}{\big|\mathcal{T}(u)\big|}\,\big\|v-u\big\|_{L^{2}(\Omega)}.
  \label{eq:ch38-defect}
\end{equation}
\index{conservation laws!guarantees for learned surrogates}
The two inequalities in \eqref{eq:ch38-defect} are the global case of a family: test the same
argument against the indicator of a control volume rather than of $\Omega$ and it bounds that
cell's defect, which is the local form an FVM forecloses by construction, as
\chapref{ch:part3-boundary-condition-enforcement-neural-oper} develops. A small aggregate
defect is therefore entirely compatible with large cancelling imbalances between neighboring
cells. The
normalization itself deserves caution: $\Delta_{\mathcal{T}}$ is meaningless for a conserved quantity
whose net value vanishes --- circulation, net charge, the integral of a vorticity anomaly --- and
must then be measured against $\int_{\Omega}|u|\,\mathrm{d}x$ instead. The surrogate compounds
the difficulty by emitting states rather than the interface fluxes
$\mathbf{F}\cdot\hat{\mathbf{n}}$ a balance would be audited against, so the cell-wise defect
cannot be read from its output without a separate flux reconstruction. Imposing the balance by
construction is the right shape of answer and this book does not supply one: the
distance-envelope device of the chapter just cited enforces a trace, not a divergence
identity. What remains is weaker and should be reported as such --- penalize the balance
residual in the objective, in the spirit of the soft penalty \eqref{eq:ch07-bc-loss}, and
release on measured $\Delta_{\mathcal{T}}$ and its cell-wise analogue alongside the acceptance metric
Implementation Pattern~\ref{imp:ch27-artifact} requires the artifact to carry.

The third obstacle is the price of the labels, which \eqref{eq:ch38-crossover} places inside
$c_{\mathrm{train}}$ rather than outside the analysis. Supervised training needs paired fields, produced by the solver the
surrogate is meant to retire, at a resolution fine enough for the budget
\eqref{eq:ch27-budget} to certify against; retraining on the same labels does not shrink the
discretization term of that budget, and storing a three-dimensional field times a long
trajectory times hundreds of realizations is often the first line item in a project plan,
ahead of the optimizer. Two levers appear in this part's own chapters. Buy labels where they
are cheap and enforce physics where they are not: coarse-data, fine-residual training
\cite{li2021physics} is intelligible precisely because \eqref{eq:ch07-apriori} says a small
weak residual plus a correct trace forces a small energy-norm error --- a separation of
statistical, discretization, and truncation contributions that \cite{subedi2024controlling}
makes explicit in a linear model. Or share an archive: the atmospheric row trains on the ERA5
reanalysis rather than on a purpose-built ensemble of solver runs, so the marginal cost of
another forecast experiment is nearly zero. The marginal epistemic cost is not, because a
reanalysis is itself the output of a data-assimilation system whose biases the surrogate
inherits with no term waiting to receive
them.\index{data generation cost}\index{carbon sequestration!uncertainty quantification over realizations}

% -----------------------------------------------------------------------------
\section{Distribution Shift and Clinical Clearance}
\label{sec:ch38-data-shift}
% -----------------------------------------------------------------------------

The obstacles so far have concerned the missing reference: a solver that would certify the
field, balance the flux, or generate the labels. The last two concern the population the
field belongs to and the institution that must accept it. The fourth is distribution shift,
the out-of-distribution (OOD) problem this part raised earlier in outline. Classical
solvers fail visibly: Newton iterations diverge, the Courant--Friedrichs--Lewy condition is
violated, the mesh tangles. A neural operator returns
a field of the right shape and the right spatial structure for an input it has never seen,
which is why
\chapref{ch:part2-production-deployment-and-high-throughput-} treats a silently
un-normalized exported artifact as the dangerous class of failure. The quantitative content is
two lines. Write $\ell(a,\theta)$ for the relative squared error of \eqref{eq:ch01-risk} at
input $a$, $\mathcal{R}_{\mu}$ for its expectation under the training measure $\mu$, and
$\mathcal{R}_{\nu}$ for the same quantity under a deployment measure $\nu$. If the loss is
bounded, with $B=\sup_{a\in\mathcal{A}}\ell(a,\theta)<\infty$, then
\begin{equation}
  \big|\mathcal{R}_{\nu}(\theta)-\mathcal{R}_{\mu}(\theta)\big|
  \;\le\;
  B\,\big|\nu-\mu\big|(\mathcal{A}),
  \label{eq:ch38-shift}
\end{equation}
where $|\nu-\mu|(\mathcal{A})$ is the total variation mass of the signed measure
$\nu-\mu$.\index{distribution shift!risk under a changed measure}\index{out-of-distribution (OOD) failure!in deployed surrogates}
Equation \eqref{eq:ch38-shift} makes the obstacle exact in three ways. It is vacuous unless
$B$ is finite, and finiteness for a relative metric requires $\|\mathcal{G}(a)\|_{\mathcal{U}}$
to stay bounded away from zero --- the qualification already attached to \eqref{eq:ch01-risk},
which a field near rest, near zero saturation, or near isoelectric potential violates. The
right-hand side belongs to the two measures and not to the network, so no further epochs on
$\mu$ shrink it; widening $\mu$ does, and that is all an augmented
training ensemble buys. And $|\nu-\mu|(\mathcal{A})$ cannot be estimated from the training
data, because that data came from $\mu$. Two responses are standard and neither is a fix in
the sense a reviewer wants. A posterior over parameters --- the Bayesian reading of a surrogate,
in the measure-based inverse-problem sense of \cite{stuart2010inverse} --- prices uncertainty in
directions the prior admits, so its calibration is a hypothesis about the prior rather than a
bound on the field. Distribution-free intervals of conformal type promise finite-sample
coverage of the form ``the next exchangeable sample falls inside with probability at least
$1-\alpha$'', the very guarantee that fails when the next sample is not exchangeable with the
calibration set, which is to say in the shifted deployment the obstacle names. Uncertainty
quantification (UQ) is best used as a tripwire inside the training distribution and not as a
certificate outside
it.\index{uncertainty quantification!for learned surrogates}\index{conformal prediction!and exchangeability}

The fifth obstacle belongs to the two rows that touch a patient. A reconstruction operator and
a cardiac wavefront surrogate enter clinical use as intended-use claims, and the evidence
reviewers expect is prospective validation across centers and heterogeneous cohorts, under
the clearance routes administered by the United States Food and Drug Administration (FDA) and
the CE marking regime for devices sold in Europe. Two features of these applications lengthen
the pathway rather than merely cost more. The reference standard moves: the full-dose
reconstruction or the fine-mesh finite element method (FEM) simulation used as truth in
evaluation is unobtainable at the point of care, so accuracy claims must be restated through
reader studies or protocol-level endpoints. And the artifact is not a fixed object: a
retrained network is a changed device, and Table~\ref{tab:synthesis_matrix} records how little
of such a change is certifiable --- two of its five cost cells print no number, none of the
three that do rests on a fully traceable pair of timings, and not one row carries a
convergence guarantee proved for the architecture as deployed; the carbon-storage cell's
\eqref{eq:ch09-jump-tail} bounds a truncation tail at a front, and is no rate theorem for
U-FNO. Neither
obstacle is peculiar to operator learning --- a conventional solver in the same indication
answers the same reviewers --- but a solver's answers are written in a vocabulary those
reviewers already accept and a surrogate's are
not.\index{regulatory clearance!medical imaging and cardiac guidance}\index{cardiac electrophysiology!ablation planning}

\takeaway{Summary of Part III}{Five domains, five architectures, three cost ratios whose
provenances the matrix caption distinguishes --- one a study-reported cost--accuracy
quotient, one formed against a member cost this Part attributes, one halved by this
Part from a study-reported pair --- two cells the record leaves empty, and one shared column of
guarantees that names, for no row, a theorem proved for the architecture that row
deploys. What operator learning demonstrably bought here is query-count independence: the
crossover $Q^{\star}$ of \eqref{eq:ch38-crossover} is cleared by the workloads that were
selected, and each field is served on the grid its query names --- under the band-limited reading
of \eqref{eq:limit_commute} that Proposition~\ref{prop:ch06-zeroshot} supplies, and not beyond
it. What it has not bought is
certification: a residual bounds an error only where a stability constant supplies
$C_{\mathrm{st}}$ in \eqref{eq:ch38-residual}, conservation holds on average rather than cell
by cell as \eqref{eq:ch38-defect} shows, label cost sits inside the amortization it
justifies, and a risk estimated under $\mu$ says nothing about $\nu$ beyond
\eqref{eq:ch38-shift}.}

% -----------------------------------------------------------------------------
\section{What the Three Parts Established, and What Remains Open}
\label{sec:ch38-closing}
% -----------------------------------------------------------------------------

\chapref{ch:part1-introduction-the-paradigm-shift-to-infinit} opened the book with a problem
rather than a method: simulation answers one question at a time, the sciences ask many, and
stretching a fixed-dimensional network over the second shape of work leaves resolution
coupling, geometry coupling, and a semantic failure in which the network has learned \emph{a}
solution instead of the rule that generates solutions from data. The remedy proposed there was
to take the solution operator \eqref{eq:ch01-solution-operator} as the object of learning and
minimize the risk \eqref{eq:ch01-risk} under a measure on the data space. That program is now
carried as far as this book can support, and the three parts' contributions are worth stating
in the order they were earned.

Part~I established when the program is possible in principle: density on compact sets for
integral-kernel operators (Theorem~\ref{thm:no_universality}), an explicit spectral rate for
the Fourier neural operator under a smoothing hypothesis (Theorem~\ref{thm:fno_convergence}),
and a precise content for mesh-independence --- the commuting diagram \eqref{eq:limit_commute},
whose zero-shot reading is confined to inputs band-limited inside the retained modes, as
Proposition~\ref{prop:ch06-zeroshot} and Figure~\ref{fig:super_resolution} record. Part~II
turned those objects into tensor programs and accounted for them: layer cost
\eqref{eq:ch8-costs} unfolded into per-field arithmetic \eqref{eq:ch27-cost}, training regimes
whose failure modes are diagnosable rather than occult, and an export artifact whose four
components state what the model was certified on. Part~III asked the question that decides
whether any of it matters, which is whether $Q$ is large enough. In aerodynamic sweeps,
atmospheric ensembles, century-scale carbon risk, electrophysiology planning, and low-dose
tomography it is, and the discretization-invariant surrogate is the mechanism that makes the
answer affordable.\index{multi-query surrogate problem!resolved by amortization}

What remains open is the boundary of the theory, marked in
\chapref{ch:part1-open-theoretical-problems-and-mathematical} and re-encountered in every row
of Table~\ref{tab:synthesis_matrix}. Sample complexity: the estimation term of
\eqref{eq:ch09-risk} and the accuracy \eqref{eq:ch09-rate} it implies are known only under
Assumption~\ref{ass:ch09-compact}, and no bound in this book is explicit in depth, channel
count, and mode count for the operators Part~III trained. Discontinuities: a saturation front
or a shock leaves the Sobolev scale that buys the rate, \eqref{eq:ch09-jump-tail} prices the
tail left behind, and \eqref{eq:ch09-gibbs} fixes the overshoot no $L^2$ number reports.
Rollouts: the accumulation in \eqref{eq:ch09-gronwall} is governed by the Lipschitz product
\eqref{eq:ch09-lipschitz}, which no loss term in Part~II controls and only the optimizer's
weight decay touches, and what a climate or
arrhythmia application reads is the weak error \eqref{eq:ch09-weak}, which can be small while
the pathwise error grows exponentially. Each roadblock of
Sections~\ref{sec:ch38-roadblocks} and~\ref{sec:ch38-data-shift} is
one of those gaps seen from the far side of a deployment: the missing stability constant is
\eqref{eq:ch38-residual} with no $C_{\mathrm{st}}$, the missing cell balance is a term the
split \eqref{eq:ch01-error-split} does not contain, and the unestimable shift is the question
\eqref{eq:ch09-risk} of Section~\ref{sec:ch09-sample-complexity} does not ask, its bound being
an $L^{2}(\mu)$ statement that says nothing about inputs outside the support of $\mu$. Nothing new is claimed here; a closing chapter that added measurements would be a sixth
application rather than a synthesis.

The reader who wants the book's components laid out rather than argued will find them in the
back matter. Appendix~\ref{app:a} collects the functional-analysis and Sobolev results Part~I
leans on, with the embeddings the convergence proofs consume. Appendix~\ref{app:b} is the
complete runnable one-dimensional Fourier neural operator for periodic heat diffusion, its
errors measured by an archived re-execution of the script, and the band-limited caveat on
resolution transfer kept visible
rather than footnoted away. Appendix~\ref{app:c} fixes the benchmark
specifications --- domain, nondimensionalization, forcing, input measure, error
metric --- against which \chapref{ch:part2-canonical-pde-benchmarks-and-reference-wor} reports,
and Appendix~\ref{app:d} consolidates the training schedules and failure modes.
Appendix~\ref{app:e} indexes the architectural variants by name, an inventory wider than
Table~\ref{tab:synthesis_matrix} and read from the architecture side instead of the domain
side. The glossary in Appendix~\ref{app:glossary} defines each technical term in its own words;
the bibliography collects every work cited in these three parts in one alphabetized sequence,
recording the venue and, wherever the publisher's record supplies one, the DOI or preprint
identifier; and the subject index carries the entries authored chapter by chapter, with
cross-references
between the theory that produced them and the applications that spent
them.\index{back matter!appendices, glossary, bibliography, index}

\vspace{1.5em}
